Software rejuvenation refers to the process of updating legacy codebases to incorporate modern programming language features while preserving their original semantics ~\cite{DBLP:conf/sofsem/PirkelbauerDS10}. Unlike general refactoring, which focuses on improving internal structure, rejuvenation addresses technical obsolescence driven by language evolution and ecosystem changes. Prior work has highlighted that such modernization efforts often involve trade-offs between adopting new language features and maintaining backward compatibility in existing systems~\cite{mendonca2026}.

Empirical studies on Python evolution have primarily focused on the adoption of specific language features. Research on type hints shows that their adoption initially emerged at low rates and is typically partial, often concentrated on function signatures rather than full codebases ~\cite{10.1145/3426422.3426981, 10.1145/3540250.3549114}. Once introduced, type annotations tend to remain stable over time, with adoption largely motivated by improved documentation and tool support such as static analysis and IDE assistance ~\cite{asarethird2026}.

%Similarly, studies on \spm have investigated its impact on code structure and readability. Results indicate that \texttt{match/case} constructs can reduce syntactic complexity compared to traditional conditional chains, although their benefits may vary depending on usage context and developer familiarity~\cite{10006855, 10589769}.

While existing studies have primarily focused on adoption rates, automated transformations, or language design, fewer works investigate the motivations behind large-scale code modernization from the developer’s perspective. Some recent efforts have begun to incorporate qualitative analysis of development artifacts to better understand intent behind changes~\cite{DBLP:conf/sbqs/CarvalhoB0M24}. However, these studies still rarely combine repository mining with direct developer feedback.

In contrast, this study investigates the motivations behind software rejuvenation in Python systems by combining repository mining with thematic analysis and triangulation of commit messages, pull requests, and developer responses. This enables a complementary perspective that connects observable code changes with the underlying intentions expressed by developers.
